\section{Fire dara description and visualization}

\subsection{Variable Definitions and Geospatial Metrics}

Understanding the physical meaning of the VIIRS variables is critical for interpreting the visualizations. The following table summarizes key columns used in the analysis:

\begin{table}[H]
    \centering
    \caption{Summary of Key VIIRS Fire Dataset Variables}\label{tab:viirs_variables}
    \begin{tabular}{llc}
    \toprule
    \textbf{Variable} & \textbf{Meaning} & \textbf{Unit} \\
    \midrule
    \texttt{bright\_ti4} & Brightness temp., Band M13 ($4\,\mu\text{m}$) & Kelvin (K) \\
    \texttt{bright\_ti5} & Brightness temp., Band M15 ($11\,\mu\text{m}$) & Kelvin (K) \\
    \texttt{frp} & Fire Radiative Power & Megawatts (MW) \\
    \texttt{scan} & Pixel size in scan (E-W) direction & Degrees \\
    \texttt{track} & Pixel size in along-track (N-S) direction & Degrees \\
    \bottomrule
    \end{tabular}
\end{table}

\subsection{Exploratory Data Analysis}
\paragraph{Feature Distribution and Correlation}

\begin{figure}[H]
    \centering
    \includegraphics[width=0.5\textwidth]{\detokenize{images/Fire/distribution of numerical columns.png}}
    \caption{Histograms illustrating the distribution of key numerical variables (\texttt{latitude}, \texttt{longitude}, \texttt{bright\_ti4}, \texttt{frp}, \texttt{scan}, \texttt{track}, \texttt{bright\_ti5}) in the merged fire dataset.}\label{fig:histograms}
\end{figure}

The distributions of the numerical columns were visualized using histograms, revealing a heavily left-skewed distribution for intensity variable \texttt{frp}, indicating that most fire detections are of low power with few extreme events, which confirms what has been shown in the previous scatter plot frp map.

\begin{figure}[H]
    \centering
    \includegraphics[width=0.5\textwidth]{\detokenize{images/Fire/heatmaps of numerical.png}}
    \caption{Heatmap of the Pearson correlation matrix for numerical variables. High correlation exists between the brightness temperatures and the pixel dimensions.}\label{fig:heatmap}
\end{figure}

A correlation matrix heatmap was generated to explore relationships between the numerical features. Key findings include:
\begin{itemize}
    \item A very high correlation (Spearman $\rho \approx 0.96$) was observed between \texttt{bright\_ti4} and \texttt{bright\_ti5}, which is expected as both measure thermal radiation from the same source, but in different bands.
    \item Positive correlations were also noted between \texttt{bright\_ti4}, \texttt{frp}, \texttt{scan}, and \texttt{track}, suggesting that larger or more intense fires (higher \texttt{frp}) often have higher brightness temperatures and are more likely to be detected in larger, more distorted pixels (higher \texttt{scan} and \texttt{track} values).
\end{itemize}
